% Conventions: unnumbered, before the numbered sections, read by every one of
% them.  The conventions task writes this file: the notation and where it
% differs from the Lean names, the letters and their one meaning each, the
% standing hypotheses (each with its \fcdefhyp in notation.sty), natural
% numbers against integers, the junk values the statements rely on and, when
% the text cites a reference work, what a reference such as \fcref{...} means
% and what its dagger marks.
\section*{Conventions}\label{sec:conventions}

Throughout, $X$ is a \fcnotion{TopologicalSpace}{topological space} and $Y$ is a
second one, and $f$ is a \fcnotion{ContinuousMap}{continuous map} from $X$ to
$Y$.  Where a distance on $X$ is needed, $X$ carries a
\fcnotion{PseudoMetricSpace}{pseudometric}, and
\fcnotion{Metric.ball}{$\ball{x}{\rho}$} is the open ball of radius $\rho$ about
a point $x$, for a real number $\rho$ greater than zero.  The letters $x$, $y$
and $z$ name points of $X$, and $x$ is the basepoint wherever a basepoint is
chosen.  The letters $x_{0}$, $x_{1}$ and $x_{2}$ name further points of $X$,
and $x'$ and $y'$ name points of $X$ tied to $x$ and $y$ by an equation.

The letters $A$ and $B$ name two subsets of $X$; $S$ names a single subset of
$X$ when only one is in play, and $T$ names a subset of $X$ containing $S$.  The
letter $I$ names an index type, $i$ and $j$ name its elements, and $U$ names a
family of subsets of $X$ indexed by $I$, whose member at $i$ is $U_{i}$.  Each
subset of $X$ carries \fcnotion{instTopologicalSpaceSubtype}{the topology it
inherits from $X$} and is a space in its own right; \fcnotion{Set.Elem}{a point
of $S$} is a point of $X$ lying in $S$, and the passage between the two is left
silent.  The map \fcnotion{ContinuousMap.subtypeVal}{$\incl{S}$} is the
inclusion of $S$ into $X$, and
\fcnotion{ContinuousMap.inclusion}{$\inclin{S}{T}$} is the inclusion of $S$ into
$T$.  The interior of $A$ is written \fcnotion{interior}{$\intr A$}.  A
\fcnotion{nhds}{neighborhood} of a point of a space is a subset of that space
whose interior contains the point, so such a subset need not be open.

The unit interval is \fcnotion{unitInterval}{$\uI$}, and the letters $r$, $s$,
$t$, $u$ and $v$ name points of it, which are called parameters.  A
\fcnotion{Path}{path} in $X$ from $x$ to $y$ is a continuous map from $\uI$ to
$X$ whose value at $0$ is $x$ and whose value at $1$ is $y$; the letters
$\gamma$ and $\delta$ name paths, a loop at $x$ is a path from $x$ to $x$, and
the \fcnotion{Set.range}{range} of a path is the set of its values.  For a path
$\gamma$ from $x$ to $y$ and a path $\delta$ from $y$ to $z$, the concatenation
\fcnotion{Path.trans}{$\gamma \ptrans \delta$} runs through $\gamma$ at twice
the speed on the first half of $\uI$ and through $\delta$ at twice the speed on
the second.  For parameters $r$ and $s$, the subpath
\fcnotion{Path.subpath}{$\subpath{\gamma}{r}{s}$} runs through $\gamma$ from
$\gamma(r)$ to $\gamma(s)$ at a constant speed.  When every value of $\gamma$
lies in a subset $S$, the path $\codres{\gamma}{S}$ is $\gamma$ read as a path
in $S$, and \fcnotion{Path.map}{$f \circ \gamma$} is $\gamma$ read as a path in
$Y$ from $f(x)$ to $f(y)$.  When $x' = x$ and $y' = y$, the path $\gamma$ is
read as a path from $x'$ to $y'$ by \fcnotion{Path.cast}{renaming} its
endpoints.

A subset of a space may be \fcnotion{IsCompact}{compact}, and it is
\fcnotion{IsPathConnected}{path connected} when one of its points is joined to
every one of its points by a path lying in it, so a path connected subset is
nonempty.  A \fcnotion{PathConnectedSpace}{path-connected space} is a nonempty
space in which any two points are joined by a path.

Two paths with the same two endpoints are homotopic when a homotopy carries one
to the other and holds both endpoints fixed throughout, and
\fcnotion{Path.Homotopic.Quotient}{$\pcls{\gamma}$} is the class of $\gamma$
under this relation; \fcnotion{Path.Homotopic.Quotient.cast}{renaming} the
endpoints of a class renames the endpoints of its representatives.  The symbol
$\ptrans$ also writes \fcnotion{Path.Homotopic.Quotient.trans}{the composition}
of two such classes, so that $\pcls{\gamma} \ptrans \pcls{\delta}$ is
$\pcls{\gamma \ptrans \delta}$.  The fundamental groupoid
\fcnotion{FundamentalGroupoid}{$\fgpd{X}$} has the points of $X$ as its objects
and the classes of paths from $x$ to $y$ as its morphisms from $x$ to $y$,
composed by concatenation; the letter $\omega$ names such a class.  The functor
\fcnotion{FundamentalGroupoid.map}{$\fgpd{f}$} sends a point to its image under
$f$ and a class to the class of its image.  The fundamental group
\fcnotion{FundamentalGroup}{$\fgrp{X}{x}$} is the group of morphisms of
$\fgpd{X}$ from $x$ to itself, the product of two of them being the second one
followed by the first, and \fcnotion{FundamentalGroup.map}{$\fgrp{f}{x}$} is the
homomorphism from $\fgrp{X}{x}$ to $\fgrp{Y}{f(x)}$ that $f$ induces.  The class
of a loop at $x$ is \fcnotion{FundamentalGroup.fromPath}{an element} of
$\fgrp{X}{x}$.  For a point $x$ lying in $S$, the group $\fgrp{S}{x}$ is the
fundamental group of $S$ based at that point of $S$.

A \fcnotion{CategoryTheory.Subgroupoid}{subgroupoid} $\mathcal{H}$ of $\fgpd{X}$
picks out, for each ordered pair of points, a set of morphisms between them, and
these sets are closed under composition and under inverses.  When $\fgpd{f}$ is
injective on objects, \fcnotion{CategoryTheory.Subgroupoid.im}{$\im \fgpd{f}$}
is the subgroupoid of the morphisms that $\fgpd{f}$ produces, and
\fcnotion{MonoidHom.range}{$\im \varphi$} is the \fcnotion{Subgroup}{subgroup}
of the values of a homomorphism $\varphi$.  Given subgroupoids $\mathcal{H}_{i}$
of $\fgpd{X}$, one for each $i$, the subgroupoid $\gen{\mathcal{H}_{i}}$ is
\fcnotion{iSup}{the smallest one containing every $\mathcal{H}_{i}$}, and the
$\mathcal{H}_{i}$ generate $\fgpd{X}$ when $\gen{\mathcal{H}_{i}}$ is
$\fgpd{X}$.  Subgroups of a group are written and generated the same way.

The letters $\mathcal{C}$ and $\mathcal{D}$ name
\fcnotion{CategoryTheory.Category}{categories}, and $\mathcal{C}$ is a
\fcnotion{CategoryTheory.Groupoid}{groupoid} where that is said, so that each of
its morphisms has an \fcnotion{CategoryTheory.inv}{inverse}.  The letter $w$
names a chosen object of $\mathcal{C}$, the letters $c$ and $c'$ name further
objects of $\mathcal{C}$, and $g$ names a morphism of $\mathcal{C}$.  The monoid
of morphisms of $\mathcal{C}$ from $w$ to itself is
\fcnotion{CategoryTheory.End}{$\End(w)$}, a group when $\mathcal{C}$ is a
groupoid, and the product of two of its elements, written with the symbol
\fcnotion{CategoryTheory.CategoryStruct.comp}{$\circ$} that composes two
morphisms of a category from right to left, is again the second one followed by
the first.  The letter $\tau$ names a family that gives, for each
object $c$, a morphism $\tau_{c}$ from $w$ to $c$.  The letter $M$ names a
\fcnotion{Monoid}{monoid}, and \fcnotion{CategoryTheory.SingleObj}{$\sobj{M}$}
is the category with one object whose morphisms from that object to itself are
the elements of $M$; one of them followed by a second composes to the product of
the second with the first.  The letters $\varphi$ and $\psi$ name
\fcnotion{MonoidHom}{homomorphisms} of monoids, and $\varphi_{A}$ and
$\varphi_{B}$ are the ones given on $A$ and on $B$.  The letter $F$ names a
\fcnotion{CategoryTheory.Functor}{functor}, $F_{i}$ the functor given on
$U_{i}$, and $F_{A}$ and $F_{B}$ those given on $A$ and on $B$.  The symbol
$\circ$ also writes \fcnotion{CategoryTheory.Functor.comp}{the composite of two
functors} and \fcnotion{MonoidHom.comp}{the composite of two homomorphisms},
read from right to left in the same way.  A square of
groups is a \fcnotion{CategoryTheory.IsPushout}{pushout} square when it is a
pushout in the \fcnotion{GrpCat}{category of groups}.

The set of natural numbers is \fcnotion{Nat}{$\N$} and starts at $0$; the
letters $n$ and $k$ name natural numbers, and a subdivision of $\uI$ into $n$
pieces is written $t_{0}, \dots, t_{n}$, its points
\fcnotion{Monotone}{nondecreasing}.  The set of integers is $\Z$ and the set of
real numbers is \fcnotion{Real}{$\R$}.  A parameter is compared with and
combined with real numbers through its value in $\R$, so that an expression such
as $2u - 1$ is read in $\R$ and no subtraction is ever taken in $\N$.  The
letter $G$ names an \fcnotion{AddCommGroup}{additive commutative group} that is
\fcnotion{LinearOrder}{linearly ordered} and
\fcnotion{IsOrderedAddMonoid}{whose addition respects that order}; $a$ and $b$
name elements of $G$ with $a$ at most $b$, $d$ names an element of $G$ used as a
step size, $nd$ is the sum of $n$ copies of $d$, and
\fcnotion{Set.Icc}{$\cint{a}{b}$} is the closed interval from $a$ to $b$.

Two constructions are given a value everywhere, including where the mathematics
they describe says nothing.  \fcnotion{Set.Icc.addNSMul}{The point reached from
$a$ by $n$ steps of size $d$} is pushed into $\cint{a}{b}$: it is $a + nd$ when
that element lies in the interval, $a$ when that element is smaller than $a$,
and $b$ when that element is larger than $b$.  The subpath
$\subpath{\gamma}{r}{s}$ is formed for every pair of parameters, so it runs
backwards along $\gamma$ when $s$ is smaller than $r$, and it is the constant
path at $\gamma(r)$ when $r$ and $s$ agree.

Four hypotheses recur, each in fixed words: that \hypCoverTwo*; that
\hypCoverFamily*; that \hypPathConnected*; and that \hypBasepoint*.  A
statement assumes one of the four only where it says so.

A number such as \fcref{Theorem 1.20} points to the result carrying that
number in the reference work.
